\chapter{Invariants, monovariants and extremal arguments}\index{invariant}

Three questions dominate olympiad problems: what survives a process
(invariant), why it must end (monovariant), and where the boundary case
lives (extremal principle). Together they turn ``show it's impossible''
into ``count the conserved quantity'' --- the cheapest proofs in the
book once you learn to look for what doesn't move.

% from-zero:begin
\begin{rulebox}[Before you start]
Chapters~2--3 and~8. \emph{Recap:} an \emph{invariant} is a quantity that never changes during a process; a \emph{monovariant} is a quantity that only moves in one direction.
\end{rulebox}
\begin{wins}[After this chapter you can]
\begin{itemize}
\item find an invariant that proves a process can never reach a target state;
\item use a monovariant to prove a process must stop;
\item apply extremal arguments (``consider the smallest\dots'') and parity colourings;
\item analyse a simple game by finding a winning strategy.
\end{itemize}
\end{wins}
\begin{principle}[How to think about this]
When you are asked whether something is possible, look for something that never changes (an invariant) or that only moves one way (a monovariant). If a target state breaks the unchanging quantity, it is impossible; if a quantity must keep shrinking, the process must stop. Try small cases until a quantity suggests itself.
\end{principle}
% from-zero:end

\section{Invariants}\index{invariant!definition}
An \emph{invariant} is a quantity (or parity, or property) unchanged by
every legal move. The discipline: state it, then check \emph{every}
move, not the moves you expect. If a target configuration has a
different value, it is unreachable --- impossibility for free.

\begin{example}[The chocolate bar]\index{chocolate bar}
Breaking an $m\times n$ bar along grid lines, one piece at a time:
prove you always need exactly $mn-1$ breaks.

\textbf{Strategy.} Don't count cuts --- count what each move preserves.
Every break splits one piece into two: the \emph{number of pieces} rises
by exactly 1 no matter where the break lands.

\textbf{Solution.} The component count is the monovariant:
$1$ piece at start, $mn$ pieces at the end, $+1$ per break ---
so $mn-1$ breaks, always (Figure~\ref{fig:choco}). Order and shape are
irrelevant; the invariant does the whole proof.
\end{example}
\begin{figure}[ht]\centering
% Figure: chocolate bar component count (3-step small multiple).
% Source for Figure~\ref{fig:choco}. TikZ only; regenerable by compiling main.tex.
\begin{tikzpicture}[
  cell/.style={draw=ink,minimum size=0.32cm,inner sep=0},
  lab/.style={font=\footnotesize\sffamily,text=slate,align=center},
]
\foreach \x in {0,1,2} \foreach \y in {0,1} \node[cell] at (0.32*\x,0.32*\y) {};
\node[lab] at (0.32,-0.5) {1 piece, 0 breaks};
\foreach \x in {0,1} \foreach \y in {0,1} \node[cell] at (2.0+0.32*\x,0.32*\y) {};
\foreach \y in {0,1} \node[cell] at (2.96,0.32*\y) {};
\node[lab] at (2.32,-0.5) {2 pieces, 1 break};
\foreach \x in {0} \foreach \y in {0,1} \node[cell] at (4.0+0.32*\x,0.32*\y) {};
\foreach \x in {1} \foreach \y in {0,1} \node[cell] at (4.64+0.32*\x,0.32*\y) {};
\node[lab] at (4.32,-0.5) {3 pieces, 2 breaks};
\node[lab] at (2.32,-1.2) {each break adds exactly one component ---\\
\#pieces $=$ \#breaks $+1$ \emph{always}};
\end{tikzpicture}

\caption{The chocolate-bar invariant. Notice you never need to know
\emph{where} the break happened --- only that components went up by 1.}
\label{fig:choco}
\end{figure}

\section{Monovariants}\index{monovariant}
A \emph{monovariant} is a measure that only moves one way --- usually a
nonnegative integer that strictly decreases. Since it can't decrease
forever, the process \emph{terminates} (Figure~\ref{fig:descent}):
termination is a corollary, not a separate proof.
\begin{figure}[ht]\centering
% Figure: monovariant descent plot (strictly decreasing nonnegative integer).
% Source for Figure~\ref{fig:descent}. TikZ only; regenerable by compiling main.tex.
\begin{tikzpicture}[
  lab/.style={font=\footnotesize\sffamily,text=slate,align=center},
  pt/.style={circle,fill=accent,inner sep=1.5pt},
]
\draw[rulegray,->] (0,0) -- (4.6,0);
\draw[rulegray,->] (0,0) -- (0,2.2);
\foreach \x/\y in {0/1.9, 0.7/1.5, 1.4/1.5, 2.1/0.9, 2.8/0.9, 3.5/0.4, 4.2/0.05}
  \node[pt] at (\x,\y) {};
\draw[ink,thick,dashed] (0,1.9) -- (0.7,1.5) -- (1.4,1.5) -- (2.1,0.9) -- (2.8,0.9) -- (3.5,0.4) -- (4.2,0.05);
\node[lab,align=center]{a nonincreasing integer measure bounded\\
below must floor out --- termination is free};
\end{tikzpicture}

\caption{A monovariant plot. Notice what makes it a proof: integer,
nonnegative, strictly dropping --- three properties, one conclusion.}
\label{fig:descent}
\end{figure}

\begin{example}[Difference reduction terminates]\index{difference reduction}
On a list of positive integers, replace two entries $a,b$ by $|a-b|$
until one number remains. Prove the process ends and the survivor
equals $\gcd$ of the list.

\textbf{Strategy.} Two measures: the \emph{sum} (strictly drops ---
monovariant, gives termination) and the \emph{gcd} (never changes ---
invariant, gives the answer).

\textbf{Solution.} Each move replaces $a+b$ by $|a-b|<a+b$: the sum is a
strictly decreasing positive integer, so the list must shrink to one
entry. Every move also preserves $\gcd$: $\gcd(a,b,\dots) =
\gcd(|a-b|,b,\dots)$. The survivor therefore equals
$\gcd(a_1,\dots,a_n)$ --- verified numerically in \S\ref{sec:invcode}.
\end{example}

\section{The extremal principle}\index{extremal principle}
On a finite set, a minimum \emph{exists}: name it and argue on it.
``Take the smallest counterexample'' is induction in extremal clothes;
``take the closest pair'' is geometry's version
(Figure~\ref{fig:extremal}). The trap: the principle needs
\emph{finiteness or well-ordering} --- ``the least $n$'' exists on
$\mathbb N$ but not on $\mathbb Z$.
\begin{figure}[ht]\centering
% Figure: extremal principle -- name the minimal witness.
% Source for Figure~\ref{fig:extremal}. TikZ only; regenerable by compiling main.tex.
\begin{tikzpicture}[
  lab/.style={font=\footnotesize\sffamily,text=slate,align=center},
  pt/.style={circle,fill=ink,inner sep=1.5pt},
]
\foreach \x/\y in {0.4/1.7, 0.9/1.1, 1.5/1.9, 1.1/0.5, 2.0/1.4, 2.5/0.8, 2.9/1.7, 3.3/0.4, 1.9/0.3}
  \node[pt] at (\x,\y) {};
\node[circle,fill=accent,inner sep=2pt] (mn) at (1.9,0.3) {};
\draw[accent,thick,dashed] (mn) -- (2.5,0.8);
\node[lab] at (2.5,-0.5) {finite set $\Rightarrow$ a minimum exists ---\\
\emph{name} the closest pair, then argue on it};
\end{tikzpicture}

\caption{The extremal principle. Notice the two-step pattern: existence
from finiteness, then every argument runs on \emph{that} witness.}
\label{fig:extremal}
\end{figure}

\begin{example}[Closest pair as a proof device]\index{closest
pair}\index{Sylvester--Gallai theorem}
Among $n\ge3$ points in the plane, not all collinear, show that some
line passes through exactly two of them (Sylvester--Gallai).

\textbf{Strategy.} Don't search all $\binom n3$ triples --- take the
minimum over all (point, line) distances and argue on \emph{that}
witness.

\textbf{Solution.} Consider all pairs $(P,L)$ where $L$ is a line
through two given points and $P$ is a given point not on $L$; such
pairs exist because the points are not all collinear, and they are
finitely many, so one pair $(P,L)$ minimizes $\operatorname{dist}(P,L)$.
Write $L=QR$ with $Q,R$ given points, and let $F$ be the foot of the
perpendicular from $P$ to $L$. Suppose $L$ contained a third given
point $S$. Then $Q,R,S$ are three points of the set on $L$, so two of
them --- call the nearer one $A$ and the farther $B$ --- lie on the
same side of $F$ (possibly $A=F$). For $X$ on $L$,
$\operatorname{dist}(X,\text{line }PB)=|XB|\sin\angle PBL$; hence
\[
\operatorname{dist}(A,PB)=|AB|\sin\angle PBL
\le|FB|\sin\angle PBL=\operatorname{dist}(F,PB)
<|PB|\sin\angle PBL=\operatorname{dist}(P,L),
\]
using $|FB|<|PB|$ (hypotenuse of the right triangle $PFB$). But $PB$
is a line through two given points and $A\notin PB$ (otherwise
$PB=L$ would contain $P$), so $(A,PB)$ is a pair with smaller
distance --- contradicting minimality. Therefore $L$ passes through
exactly two of the points.

\textbf{Significance.} The witness does the work; the extremal
principle supplied it. Note what the minimum had to be \emph{of}: the
closest \emph{pair} alone is not enough (three collinear points defeat
the naive version), so the quantity minimized is upgraded to a
point--line distance.
\end{example}

\section{Advanced pigeonhole}\index{pigeonhole!generalized}
Generalized form: if $N$ objects land in $k$ boxes, some box holds
$\ge\lceil N/k\rceil$. Averaging form: if numbers average $\bar x$,
some entry is $\ge\bar x$ (and some $\le$). These are the same
statement; olympiad use is always of the form ``so one object
\emph{must} have the property''.

\section{Coloring and parity arguments}\index{coloring argument}
Coloring turns structure into arithmetic: assign labels so every legal
move has a known signature, then count.
\begin{example}[The domino impossibility]\index{domino tiling}
A $10\times10$ board minus two \emph{opposite} corners: can dominoes
tile it?

\textbf{Strategy.} Checkerboard it: every domino covers one square of
each color, so a tiling requires equal color counts.

\textbf{Solution.} Opposite corners share a color (even side length);
the board has $48$ of one color and $50$ of the other --- no tiling
exists (Figure~\ref{fig:checker}). The proof is four words long once the
coloring is found: \emph{count the colors}.
\end{example}
\begin{figure}[ht]\centering
% Figure: checkerboard with two opposite corners removed.
% Source for Figure~\ref{fig:checker}. TikZ only; regenerable by compiling main.tex.
\begin{tikzpicture}[
  lab/.style={font=\footnotesize\sffamily,text=slate,align=center},
]
\foreach \i in {0,...,4} \foreach \j in {0,...,4}
  \pgfmathsetmacro{\c}{mod(\i+\j,2) ? "rulegray!45" : "white"}
  \fill[\c] (0.3*\i,0.3*\j) rectangle (0.3*\i+0.3,0.3*\j+0.3);
\draw[ink] (0,0) rectangle (1.5,1.5);
\fill[white] (0,0) rectangle (0.3,0.3);
\fill[white] (1.2,1.2) rectangle (1.5,1.5);
\draw[accent,thick] (0,0) rectangle (0.3,0.3);
\draw[accent,thick] (1.2,1.2) rectangle (1.5,1.5);
\node[lab] at (0.75,-0.5) {removed corners share a color:\\
$24$ vs $26$ squares --- no domino tiling possible};
\end{tikzpicture}

\caption{The coloring argument. Notice the impossibility needs only the
imbalance --- you never exhibit a failed tiling.}
\label{fig:checker}
\end{figure}

\section{Games}\index{games!winning positions}
Positions split into \emph{winning} (a move to a losing position
exists) and \emph{losing} (every move lands on a winning position) ---
compute backwards from terminal states (Figure~\ref{fig:nim}).
\emph{Strategy stealing} and \emph{pairing} handle existence claims:
mirror the opponent's moves across a center (Figure~\ref{fig:pair}) so
their play produces yours.
\begin{figure}[ht]\centering
% Figure: pairing strategy -- mirrored cells across center.
% Source for Figure~\ref{fig:pair}. TikZ only; regenerable by compiling main.tex.
\begin{tikzpicture}[
  lab/.style={font=\footnotesize\sffamily,text=slate,align=center},
]
\draw[ink] (0,0) grid[step=0.35] (2.1,1.05);
\fill[accent!40] (0.35,0.7) rectangle (0.7,1.05);
\fill[accent!40] (1.4,0) rectangle (1.75,0.35);
\draw[accent,thick,<->] (0.52,0.87) -- (1.57,0.17);
\node[circle,fill=ink,inner sep=1pt] at (1.05,0.52) {};
\node[lab] at (1.05,-0.5) {mirror each move through the center:\\
the second player's reply is already on the board};
\end{tikzpicture}

\caption{The pairing strategy. Notice whose move it is: the second
player's reply exists because the board was symmetric before the first
player touched it.}
\label{fig:pair}
\end{figure}
\begin{figure}[ht]\centering
% Figure: take-away game tree (take 1 or 2 from n=5).
% Source for Figure~\ref{fig:nim}. TikZ only; regenerable by compiling main.tex.
\begin{tikzpicture}[
  nd/.style={draw=ink,circle,inner sep=1.5pt,font=\footnotesize\sffamily},
  lab/.style={font=\footnotesize\sffamily,text=slate,align=center},
  w/.style={circle,draw=accent,inner sep=1.5pt,font=\footnotesize\sffamily\bfseries},
]
\node[nd] (r) at (0,0) {5};
\node[w] (a) at (-1.3,-0.9) {4};
\node[w] (b) at (1.3,-0.9) {3};
\node[nd] (c) at (-1.9,-1.9) {3};
\node[nd] (d) at (-0.7,-1.9) {2};
\node[nd] (e) at (0.7,-1.9) {2};
\node[nd] (f) at (1.9,-1.9) {1};
\draw[rulegray] (r) -- (a) (r) -- (b) (a) -- (c) (a) -- (d) (b) -- (e) (b) -- (f);
\node[lab,align=center]{take 1 or 2: $n$ is losing iff $3 \mid n$ ---\\
from $n{=}5$ every move hands the opponent a multiple};
\end{tikzpicture}

\caption{Take-away game tree ($n=5$, take 1 or 2). Notice the period:
losing positions repeat mod 3 --- the invariant-style pattern inside
every impartial game.}
\label{fig:nim}
\end{figure}
\begin{example}[Pairing the second player]\index{pairing strategy}
Two players alternately place identical coins on a round table without
overlap; the last to place wins. Who wins?

\textbf{Strategy.} Don't design a first-player attack --- hand the
second player a mirror.

\textbf{Solution.} First player places a coin at the center; thereafter
whatever the opponent plays is answered by the centrally symmetric
spot. Every opponent move creates a legal mirror move, so the first
player always has a reply and wins. The symmetry does everything after
move one.
\end{example}

\section{Summary}
\begin{center}\small
\begin{tabular}{@{}ll@{}}
\toprule
\textbf{Tool} & \textbf{Kills which question} \\
\midrule
Invariant & ``can it reach?'' --- different value, unreachable \\
Monovariant & ``does it end?'' --- strictly drops, must floor \\
Extremal & ``does it exist?'' --- name the minimum \\
Pigeonhole & ``must one have P?'' --- $\lceil N/k\rceil$ \\
Coloring & ``can it tile/balance?'' --- count the colors \\
W/L positions & ``who wins?'' --- backtrack from terminal \\
Pairing/mirror & ``who wins?'' --- borrow their move \\
\bottomrule
\end{tabular}
\end{center}

\section{Code: watching the invariant}\label{sec:invcode}\index{invariant\_sim}
\texttt{invariant\_sim.py} runs the difference-reduction process and
prints the list, the dropping sum (monovariant) and the constant gcd
(invariant) at each step.
\begin{checkpoint}
Predict the survivor before running: $[12,18,30,42]$ should leave its
gcd. Then run \texttt{python3 invariant\_sim.py}.
\end{checkpoint}
\begin{lstlisting}
start: [12, 18, 30, 42]
t=1: [12, 18, 12]  sum=42  gcd=6 (invariant 6)
t=2: [6, 12]  sum=18  gcd=6 (invariant 6)
t=3: [6]  sum=6  gcd=6 (invariant 6)
survivor = 6 == gcd([12,18,30,42]) = 6: True
\end{lstlisting}
The sum column is the termination certificate; the gcd column is the
answer. The test script \texttt{test\_16\_inv.sh} re-runs the sim and
greps the survivor line.

\section{Common mistakes}\index{invariant!mistakes}
\begin{itemize}
\item \emph{Checking the moves you meant, not all legal moves.} An
invariant must survive \emph{every} legal transition --- one unchecked
case voids it.
\item \emph{Extremal on an infinite set.} ``Take the least counterexample''
needs well-ordering: valid on $\mathbb N$, meaningless on $\mathbb R$.
\item \emph{Draw $\Rightarrow$ strategy.} Showing a position isn't losing
for player one doesn't hand player two a strategy --- who moves first
is part of the claim.
\end{itemize}

\section*{Exercises}
\subsection*{Inline}
\begin{exercise}[easy]\label{exr:iv1}
A stack of coins: each move splits one stack into two. Prove finishing
takes $n-1$ moves for $n$ single coins.
\end{exercise}
\begin{solution}
Same invariant as the chocolate bar: components $+1$ per move.
\end{solution}
\begin{exercise}[easy]\label{exr:iv2}
On a chessboard, a rook moves any number of squares. Is ``color of the
square'' an invariant under a rook move?
\end{exercise}
\begin{solution}
No --- a rook move may end on either color, so color is not invariant
under rook moves. (A knight move, by contrast, always flips color;
a quantity that changes predictably is a \emph{monovariant}, the next
section's tool.)
\end{solution}
\begin{exercise}[medium]\label{exr:iv3}
Start with $\{1,2,\dots,2n\}$ on a board; each move erases two numbers
$a,b$ and writes $|a-b|$. Prove the last number is even iff $n$ is even.
\end{exercise}
\begin{solution}
Parity of the sum is invariant (sum drops by $2\min$); initial sum
$n(2n+1)$ has parity of $n$.
\end{solution}
\begin{exercise}[medium]\label{exr:iv4}
Show the take-away game (take 1 or 2) is losing exactly at $3\mid n$.
\end{exercise}
\begin{solution}
From $3k$ every move hits $3k\pm\{1,2\}$; from a non-multiple, take the
remainder to land on a multiple. Induction on $n$.
\end{solution}
\subsection*{Additional}
\begin{exercise}[easy]\label{exr:ia1}
Prove an $8\times8$ board minus two \emph{adjacent} corners \emph{can}
be tiled by dominoes.
\end{exercise}
\begin{solution}
Adjacent corners differ in color: $31$--$31$ balance; exhibit a snake
tiling row by row.
\end{solution}
\begin{exercise}[easy]\label{exr:ia2}
Is $a+b$ invariant under replacing $a,b$ by $a-b$? Check every move
type.
\end{exercise}
\begin{solution}
No: $a+b \to (a-b)+b = a$ changes it; the sum \emph{decreases} ---
that's the monovariant, not an invariant.
\end{solution}
\begin{exercise}[easy]\label{exr:ia3}
Among $13$ integers, two are congruent mod 12 --- state the boxes.
\end{exercise}
\begin{solution}
Residues mod 12 are the 12 boxes; 13 pigeons.
\end{solution}
\begin{exercise}[easy]\label{exr:ia4}
Give a nonnegative integer monovariant for the Euclidean algorithm.
\end{exercise}
\begin{solution}
The smaller argument $b$ strictly drops each division step; it's a
positive integer, so Euclid terminates.
\end{solution}
\begin{exercise}[medium]\label{exr:ia5}
Prove every $2^n\times2^n$ board minus one square is tileable by
L-trominoes (invariant/induction form).
\end{exercise}
\begin{solution}
Induction: quarter the board, put one L at the center covering one
square of each of the three ``complete'' quarters; each quarter is now
a smaller instance.
\end{solution}
\begin{exercise}[medium]\label{exr:ia6}
At a party any two people have a common friend; prove someone knows
everyone (extremal hint).
\end{exercise}
\begin{solution}
Take a person $M$ of maximum degree; if some $X$ weren't $M$'s friend,
$M$ and $X$ share a friend of degree $> \deg M$.
\end{solution}
\begin{exercise}[medium]\label{exr:ia7}
$51$ numbers from $\{1,\dots,100\}$: show two are consecutive ---
identify the coloring/box structure.
\end{exercise}
\begin{solution}
Pairs $\{1,2\},\{3,4\},\dots,\{99,100\}$ are 50 boxes; 51 picks force a
full pair.
\end{solution}
\begin{exercise}[medium]\label{exr:ia8}
Winning strategy for one-pile Nim: take any number up to $k$. State and
prove the losing set.
\end{exercise}
\begin{solution}
Losing iff $(k+1)\mid n$: multiples map to non-multiples and back,
same induction as take-1-or-2.
\end{solution}
\begin{exercise}[medium]\label{exr:ia9}
Implement-or-derive: extend \texttt{invariant\_sim.py} so the move
picks a random pair each step (seeded), and confirm the survivor is
still the gcd.
\end{exercise}
\begin{solution}
The invariant doesn't care which pair moves; random or largest-first,
the survivor is $6$.
\end{solution}
\begin{exercise}[hard]\label{exr:ia10}
Prove the circle-difference process (replace each $a_i$ by
$|a_i - a_{i+1}|$ on an even cycle of integers) reaches all zeros
when $n$ is a power of two, via a monovariant + invariant combo.
\end{exercise}
\begin{solution}
Max is nonincreasing; parity analysis (all entries become even after
$\le$ a few steps, then factor out powers of two) forces max to hit 0.
The invariant is gcd-divisibility; the monovariant is max after
rescaling.
\end{solution}
\begin{exercise}[hard]\label{exr:ia11}
Strategy stealing: prove the first player in tic-tac-toe (on an
infinite board, get 5 in a row) cannot be worse off than the second.
\end{exercise}
\begin{solution}
Suppose second player had a winning strategy; first plays an arbitrary
opening move, then ``becomes'' the second player --- an extra mark can
never hurt in a strong positional game.
\end{solution}
\begin{exercise}[hard]\label{exr:ia12}
Prove there is no perfect cuboid argument's first step: in a set of
distinct positive integers whose pairwise sums are all squares, an
extremal argument bounds the set --- give the extremal witness to name.
\end{exercise}
\begin{solution}
Take the largest element $M$: every partner's sum lies in
$(M, 2M)$, so partner squares are restricted to squares in that
interval --- finitely many candidates force a bound. (Open-problem
flavor; the point is naming the witness.)
\end{solution}
\begin{exercise}[easy]\label{exr:i13}
A chessboard has two opposite corners removed. Show 31 dominoes cannot
tile it, and name the invariant.
\end{exercise}
\begin{solution}
Each domino covers one black and one white square; the removed corners
share a colour, leaving $30$ of one colour vs $32$ of the other ---
the colour-difference invariant $|W-B|=2$ survives every placement and
no tiling can reach $0$.
\end{solution}
\begin{exercise}[medium]\label{exr:i14}
Numbers $1,\dots,10$ sit on a circle. A move picks two adjacent
numbers and swaps them. Prove no sequence of moves arranges them in
decreasing order clockwise if they start increasing --- or find a
counterexample and diagnose what the invariant actually is.
\end{exercise}
\begin{solution}
Adjacent transpositions generate $S_n$, so decreasing order
\emph{is} reachable --- the puzzle as stated is a trick with a
counterexample, not an impossibility. Where parity \emph{does} supply
a real invariant is the 15-puzzle: the permutation parity combined with
the blank's position is preserved, and it blocks exactly half of all
configurations. Lesson: an impossibility claim needs an invariant that
genuinely never changes --- guess first, verify the invariant before
believing it.
\end{solution}
\begin{exercise}[hard]\label{exr:i15}
Show the sum of the numbers on a $n$-nim position (in binary XOR,
not decimal) is the true invariant/monovariant: prove a position is
winning iff the nim-sum is nonzero, from the definitions.
\end{exercise}
\begin{solution}
Terminal position has nim-sum $0$; from any nonzero position a single
heap reduction makes the XOR zero (pick the heap containing the top
bit); from a zero position every move breaks a bit balance and leaves
nonzero. Winning positions are exactly the nonzero ones.
\end{solution}

\section*{Notes}
Invariants, colorings, the extremal principle and game positions follow
Engel~\cite{engel} and Zeitz~\cite{zeitz}; the state-machine
invariant framing is 6.042J's~\cite{mcs}; pairing and
strategy-stealing examples are Cut-the-Knot's games
index~\cite{cuttheknot}. Contest uses: 2024 AIME I P3; 2023 AIME I P14;
2021 AIME I P10; 2024 AMC 12B P5; 2024 AMC 12A P22; 2010 AIME I P12;
2022 AIME I P9; 2025 AIME II P3; 2023 AIME II P11; USAJMO 2024 P2, P4;
2023 USAJMO P3, P4, P5.

\section*{Further reading}
\begin{itemize}
\item A.~Engel, \emph{Problem-Solving Strategies}, invariants and
coloring chapters --- the canonical drill source for this chapter.
\item P.~Zeitz, \emph{The Art and Craft of Problem Solving} ---
invariants and the extremal principle with worked strategy.
\item E.~Lehman, F.~T.~Leighton, A.~R.~Meyer, \emph{Mathematics for
Computer Science}, state-machine invariants chapter --- free PDF; the
CS framing of invariants:
\url{https://courses.csail.mit.edu/6.042/spring18/mcs.pdf}
\item MIT OpenCourseWare, \emph{6.042J} ---
\url{https://ocw.mit.edu/6-042JF10} --- the matching lectures.
\item Y.~Zhao's olympiad handouts ---
\url{https://yufeizhao.com/olympiad/} --- invariant and extremal
problem sets.
\item Cut-the-Knot's games index ---
\url{https://www.cut-the-knot.org} --- pairing and strategy-stealing
examples in interactive form.
\end{itemize}

\begin{rulebox}[Watch instead]
\begin{itemize}
\item \textbf{ChessNetwork --- ``Mutilated chessboard problem''},
4:07, \url{https://www.youtube.com/watch?v=bS3EqMcIlAQ} --- the
colouring invariant of the Coloring-arguments section in its purest
form. (Verified 2026-10-11.)
\item \textbf{Mathologer --- ``The Arctic Circle Theorem''}, 51:49,
\url{https://www.youtube.com/watch?v=Yy7Q8IWNfHM} --- opens on the same
mutilated chessboard, then shows how far tiling invariants go.
(Verified 2026-10-11.)
\end{itemize}
\end{rulebox}

\chapternotes
